% Generated draft: requires mathematical and source review.
\documentclass[11pt]{article}
\usepackage[margin=25mm]{geometry}
\usepackage{fontspec}
\setmainfont{Latin Modern Roman}
\usepackage{amsmath,amssymb,mathtools,unicode-math}
\setmathfont{Latin Modern Math}
\usepackage{enumitem,needspace,xurl,hyperref}
\setlength{\parindent}{0pt}
\setlength{\parskip}{5pt}
\setlength{\emergencystretch}{3em}
\newcommand{\N}{\mathbb{N}}
\newcommand{\Z}{\mathbb{Z}}
\newcommand{\Q}{\mathbb{Q}}
\newcommand{\R}{\mathbb{R}}
\newcommand{\C}{\mathbb{C}}
\newcommand{\F}{\mathbb{F}}
\newcommand{\A}{\mathbb{A}}
\newcommand{\PP}{\mathbb{P}}
\newcommand{\E}{\mathbb{E}}
\newcommand{\eps}{\varepsilon}
\newcommand{\dd}{\,\mathrm{d}}
\begin{document}
\section*{UnsolvedMath ID: 20000538}
\textbf{Source problem code:} catalogue entry
\subsection{Definitions and mathematical statement}
Camil Muscalu's AIM question asks for a multilinear extension of Stein--Wainger polynomial Carleson theory. In the stated bilinear example let \(m=\widehat K\) satisfy
\[
|\partial^\alpha m(\xi)|\leq C_\alpha|\xi|^{-|\alpha|}
\quad(\xi\in\R^2\setminus\{0\})
\]
for sufficiently many multi-indices. For Schwartz functions \(f,g\) on \(\R\), the target is
\[
\left\|\sup_{\lambda\in\R}\left|
\operatorname{p.v.}\iint f(x-t)g(x-s)K(t,s)e^{i\lambda t^2s^2}\,dt\,ds
\right|\right\|_{L^2(\R)}
\lesssim\|f\|_{L^4}\|g\|_{L^4}.
\]
The principal-value/truncation convention and number of derivative bounds are not specified in the source; they need to match the intended singular kernel. Its displayed Li--Muscalu background formula, reference [11], loses the Fourier synthesis factor and has a three-argument expression for a two-variable multiplier. The natural two-parameter modulation form would use \(m(\xi_1-N_1,\xi_2-N_2)\) and Fourier synthesis in \(x\), but this is an explicit interpretation, not a verified verbatim formula. The source attributes an \(L^4\times L^4\to L^2\) bound for that background maximal operator to Li--Muscalu.

For \(F=f\otimes g\), define the ambient convolution family
\[
A_\lambda F(x_1,x_2)=\operatorname{p.v.}\iint
F(x_1-t,x_2-s)K(t,s)e^{i\lambda t^2s^2}\,dt\,ds.
\]
The target integrand is exactly \(A_\lambda F(x,x)\). The report proposes a conditional diagonal-trace reduction: assuming the stated ambient \(L^2(\R^2)\) degree-four polynomial maximal estimate, for every \(\sigma>1/2\) the diagonal maximal expression is bounded by
\[
C_\sigma\|F\|_{X^\sigma},\qquad
\|F\|_{X^\sigma}^2=
\iint(1+|\xi-\eta|^2)^\sigma|\widehat F(\xi,\eta)|^2\,d\xi\,d\eta.
\]
For tensors, \(\widehat F=\widehat f\,\widehat g\). This normal-direction Sobolev input is stronger/different than the requested \(L^4\times L^4\) data; it does not prove the original estimate. The endpoint \(\sigma=1/2\) is not claimed. The report also states the same trace principle for a suitably defined translation-invariant maximal convolution family assumed \(L^2\)-bounded; complete hypotheses are not printed.

Independently, it claims the desired bilinear bound for each dyadic multiplier annulus, uniformly over its scale and all \(\lambda\). For finite collections or \(\ell^1\)-summable coefficients, the associated bound is the sum of the annular bounds, with the corresponding finite count or coefficient \(\ell^1\)-norm. No summation uniform over an unrestricted infinite collection, and no unrestricted original-scale \(L^4\times L^4\) theorem, is asserted.
\subsection{Short English statement}
Seek an \(L^4\times L^4\to L^2\) polynomially modulated bilinear maximal bound. The report gives a conditional ambient-to-diagonal Sobolev estimate and scale-local bounds, while leaving the unrestricted input-norm estimate open.
\subsection{Sources}
\begin{itemize}
\item AIM workshop ``Carleson theorems and multilinear operators,'' 2015, Question 11, attributed to Camil Muscalu: \url{https://aimath.org/pastworkshops/multilinopsproblems.pdf}. No PDF was accessed here.
\item Local canonical registry \texttt{aim-analysis-notes.json}, \texttt{notes[69]}, tag \texttt{question}; permanent UnsolvedMath ID \(20000538\), source code \texttt{catalogue entry}. Original equations are labelled (0.13)--(0.16) and cite Li--Muscalu [11] as background. Superscripts, inequality symbols, multiplier arguments, and Fourier synthesis are damaged; the clean statement is absent. No full entry for [11] is supplied.
\item Supplied research attempt 1 proposes the conditional diagonal-trace and annular reductions and invokes an ambient degree-four polynomial maximal theorem without a full authoritative citation. It is machine generated, not peer reviewed, and requires independent verification before citation.
\end{itemize}
\end{document}
